%%%%%%%%%%%%%%%%%%%%%%%%%%%%%%%%%%%%%%%%%%
\section{Conclusions and Future Work}
\label{sec:conclusion}

This study examined how linguistic granularity and rule-base structure relate to the mathematical properties and optimization behavior of static Mamdani fuzzy controllers integrated into particle swarm optimization. By combining a formal characterization of complementary rule-base structures with a controlled experimental evaluation across the Classical F1--F23 benchmark suite, the work provides mathematical and empirical evidence to inform the design of fuzzy inertia-control mechanisms for bio-inspired optimization.

The global ranking places FCS-3-R1 first (mean Friedman rank 4.77) and classical PSO second (6.36); both fall within the same critical-difference group. At the instance level, the best FCS configuration ranks ahead of PSO on 50 of 62 cases, but this is an oracle comparison because the selected configuration changes after the instance is known. The fixed-subset leaders differ: FCS-3-R5 on scalable $S_1$ and FCS-5-R6 on fixed-dimensional $S_2$. Observed rank concordance is weak ($\tau=0.095$, $p=0.451$); the applied test does not establish a positive association, and its non-significant result does not establish independence. These results describe heterogeneity within this benchmark but do not establish or rule out transferability to other problem sets.

FCS-3-R1 is nominally diversity-agnostic in its consequent table, but the implemented min--max Mamdani inference leaves its effective output weakly diversity-dependent. Its leading rank and near-0.5 mean inertia therefore do not demonstrate that diversity feedback improves PSO, nor do they imply a trajectory identical to the classical linear schedule. Within the evaluated configurations, adding explicit diversity trends does not invariably improve ranking; this is not evidence that diversity feedback is unnecessary in general. Proposition~\ref{prop:mirror} characterizes complementary mirror-pair surfaces; the Holm-adjusted independent-sample comparisons show that this structural symmetry does not ensure equivalent optimization performance.

For the evaluated fuzzy controllers and benchmark functions, increasing linguistic granularity did not consistently improve optimization performance or reduce performance differences among rule-base patterns. These findings indicate that greater linguistic resolution, despite increasing the number of rules, does not by itself ensure a more effective inertia-control policy. The organization of the consequent rules is therefore an important consideration when designing fuzzy controllers for bio-inspired optimization. Together with the observed differences across benchmark groups, the results suggest that controller design should consider both the intended inertia response and the characteristics of the optimization problem. These conclusions remain descriptive and specific to the evaluated configurations; they do not establish independent causal effects, statistical equivalence, or a formal granularity-by-rule interaction.


The observed differences among fixed controller configurations provide a basis for investigating dynamic configuration-selection mechanisms in future fuzzy-controlled bio-inspired optimization systems. Such mechanisms could use information available during the search to select among predefined rule-base structures or linguistic granularities according to the evolving optimization state. Unlike the controllers evaluated in this study, whose rule bases and linguistic partitions remain fixed throughout each run, these future approaches would introduce structural changes during optimization and would require independent experimental validation. Further research should also examine the granularity-by-rule interaction through a predefined inferential design, specifying the estimand, experimental unit, and appropriate treatment of repeated measurements. Validation on independent benchmark suites and applied optimization problems would be necessary to assess whether the observed structural relationships extend beyond the classical F1--F23 functions.

